
\documentclass{article}
\usepackage[utf8]{inputenc}
\usepackage[spanish]{babel}
\usepackage{geometry}
\usepackage{booktabs}
\usepackage{xcolor}
\geometry{a4paper, margin=1in}

\title{\color{red}Reporte de Entrenamiento 5 Horas \\ \huge AG-RACE}
\author{Antigravity AI Assistant}
\date{\today}

\begin{document}

\maketitle

\section{Configuración de la Sesión}
Entrenamiento con SepCMA sobre política compacta, evaluación multi-seed fija por generación y
penalización por varianza inter-seed para mejorar estabilidad.

\begin{itemize}
    \item \textbf{Unidades en pista:} 5 Agentes Neuronales + 1 Baseline.
    \item \textbf{Algoritmo:} SepCMA (variante diagonal de CMA-ES).
    \item \textbf{Seeds por evaluación:} 5
    \item \textbf{Pasos por episodio:} 400
    \item \textbf{Población:} 8
\end{itemize}

\section{Métricas Finales}
\begin{itemize}
    \item \textbf{Generaciones totales:} 42
    \item \textbf{Tiempo total:} 600.89 minutos
    \item \textbf{Mejor fitness validado:} -5327.09
    \item \textbf{Parámetros de la política:} 13374
\end{itemize}

\section{Histórico reciente}
\begin{table}[h!]
\centering
\begin{tabular}{@{}ccccccc@{}}
\toprule
\textbf{Gen} & \textbf{Fit} & \textbf{IA Avg} & \textbf{IA Std} & \textbf{Baseline} & \textbf{Gap} & \textbf{Sigma} \\ \midrule
12 & -8969.1 & -13.9 & 398.4 & 59156.5 & -59170.4 & 0.1991 \\
13 & -14264.1 & 174.7 & 428.0 & 95862.6 & -95687.9 & 0.1991 \\
14 & -12598.5 & 369.7 & 610.9 & 86010.0 & -85640.3 & 0.1990 \\
15 & -16138.9 & 627.9 & 506.4 & 111731.6 & -111103.6 & 0.1989 \\
16 & -14378.9 & 743.3 & 567.5 & 100801.4 & -100058.1 & 0.1989 \\
17 & -7596.8 & 1513.2 & 701.8 & 61310.8 & -59797.6 & 0.1988 \\
18 & -13145.1 & 416.5 & 565.2 & 90073.4 & -89656.9 & 0.1987 \\
19 & -12078.2 & 867.4 & 478.6 & 86533.4 & -85666.0 & 0.1987 \\
20 & -8908.7 & 915.1 & 588.0 & 65623.5 & -64708.3 & 0.1986 \\
21 & -8375.1 & 1448.3 & 1073.0 & 65506.8 & -64058.5 & 0.1986 \\
22 & -10545.1 & 898.7 & 716.7 & 76234.6 & -75335.9 & 0.1985 \\
23 & -11804.2 & -35.9 & 356.7 & 77943.1 & -77979.1 & 0.1984 \\
24 & -12123.0 & 261.7 & 423.0 & 82262.9 & -82001.2 & 0.1984 \\
25 & -8266.3 & 735.6 & 619.0 & 59922.6 & -59187.0 & 0.1983 \\
26 & -12886.7 & 731.4 & 513.7 & 90833.7 & -90102.3 & 0.1982 \\
27 & -7314.2 & 1564.7 & 1104.2 & 59284.5 & -57719.9 & 0.1982 \\
28 & -9168.2 & 1560.2 & 854.0 & 71944.1 & -70384.0 & 0.1981 \\
29 & -8971.7 & 1497.1 & 838.3 & 70171.0 & -68673.9 & 0.1981 \\
30 & -8184.0 & 1166.3 & 699.6 & 62568.9 & -61402.7 & 0.1980 \\
31 & -10572.1 & 1918.7 & 1116.5 & 83702.7 & -81783.9 & 0.1979 \\
32 & -9407.5 & 2002.6 & 856.9 & 76927.1 & -74924.5 & 0.1979 \\
33 & -8632.0 & 1353.4 & 1091.7 & 66467.4 & -65114.0 & 0.1978 \\
34 & -11151.1 & 1475.5 & 536.6 & 84937.7 & -83462.1 & 0.1978 \\
35 & -10212.5 & 1699.9 & 1131.1 & 79608.3 & -77908.3 & 0.1977 \\
36 & -9053.5 & 1442.9 & 838.3 & 70301.2 & -68858.3 & 0.1976 \\
37 & -10020.7 & 1335.9 & 672.9 & 76149.4 & -74813.5 & 0.1976 \\
38 & -9423.5 & 1395.3 & 814.3 & 72434.7 & -71039.4 & 0.1975 \\
39 & -10566.1 & 1862.1 & 877.4 & 83546.9 & -81684.9 & 0.1975 \\
40 & -15391.0 & 1554.3 & 824.2 & 113423.5 & -111869.3 & 0.1974 \\
41 & -16639.2 & 1730.0 & 739.9 & 123204.9 & -121474.9 & 0.1973 \\

\bottomrule
\end{tabular}
\end{table}

\section{Lectura}
La meta no es maximizar un pico aislada, sino mejorar retorno medio y reducir sensibilidad entre
semillas. Por eso se reporta la desviación estándar inter-seed además del promedio.

\end{document}
